\documentclass[a4paper]{article}
% generated by Docutils <http://docutils.sourceforge.net/>
\usepackage{cmap} % fix search and cut-and-paste in Acrobat
\usepackage{ifthen}
\usepackage[T1]{fontenc}
\usepackage[utf8]{inputenc}
\setcounter{secnumdepth}{0}
\usepackage{tabularx}

%%% Custom LaTeX preamble
% PDF Standard Fonts
\usepackage{mathptmx} % Times
\usepackage[scaled=.90]{helvet}
\usepackage{courier}

%%% User specified packages and stylesheets
\input{latex2e.tex}

%%% Fallback definitions for Docutils-specific commands

% Provide a length variable and set default, if it is new
\providecommand*{\DUprovidelength}[2]{
  \ifthenelse{\isundefined{#1}}{\newlength{#1}\setlength{#1}{#2}}{}
}

% width of docinfo table
\DUprovidelength{\DUdocinfowidth}{0.9\linewidth}

% field list environment (for separate configuration of `field lists`)
\ifthenelse{\isundefined{\DUfieldlist}}{
  \newenvironment{DUfieldlist}%
    {\quote\description}
    {\enddescription\endquote}
}{}

% hyperlinks:
\ifthenelse{\isundefined{\hypersetup}}{
  \usepackage[colorlinks=true,linkcolor=blue,urlcolor=blue]{hyperref}
  \usepackage{bookmark}
  \urlstyle{same} % normal text font (alternatives: tt, rm, sf)
}{}
\hypersetup{
  pdftitle={Curriculum Vitae of Giulio Bottazzi},
  pdfauthor={Giulio Bottazzi}
}

%%% Body
\begin{document}
\title{Curriculum Vitae of Giulio Bottazzi%
  \label{curriculum-vitae-of-giulio-bottazzi}}
\author{}
\date{}
\maketitle

% Docinfo
\begin{center}
\begin{tabularx}{\DUdocinfowidth}{lX}
\textbf{Author}: &
	Giulio Bottazzi \\
\textbf{Contact}: &
	<\href{mailto:g.bottazzi@santannapisa.it}{g.bottazzi@santannapisa.it}> \\
\textbf{Date}: &
	29 February 2020 \\
\end{tabularx}
\end{center}

% -*- Mode: rst -*-

\textbf{Date of Birth:} 30 January 1969

\textbf{Nationality:}   Italian


\section{Brief BIO%
  \label{brief-bio}%
}

Giulio Bottazzi is a Professor of Economic Policy in the Social Science
Department of the Scuola Superiore Sant'Anna in Pisa. He graduated from
Parma University and holds a PhD in Theoretical Physics from Milano
University. He is the author of several papers published in
international journals. His research interests range from industrial
dynamics to financial economics, encompassing experimental economics
(coordination and bounded rationality), economic geography (industrial
districts and localized spillovers), theoretical and empirical aspects
of market interactions (firm growth and diversification; competition and
technical efficiency) and applied econometrics (characterization of
extreme events). He has been the Principal Investigator and
Co-Investigator of projects financed by the Italian Ministry of
University and Research, the European Commission and the Tuscany
Regional Council, both in partnership with other research institutions
and private companies.


\section{Academic Career%
  \label{academic-career}%
}

\begin{description}
\item[{\emph{November 2013-October 2019} Dean of the Academic Class of Social}] 
Sciences, Scuola Superiore Sant'Anna, Pisa, Italy

\end{description}

\emph{since June 2013} Full Professor of Economic Policy, Scuola Superiore
Sant'Anna, Pisa, Italy

\emph{2010-2013} Deputy director of the Master Program in Economics
jointly offered by Scuola Sant'Anna and Universita' di Pisa.

\emph{since 2008} Research associate, BRICK Center, Collegio Carlo Alberto,
Torino, Italy.

\emph{2007-2012} Invited Professor, Université Pantheon-Assas Paris II.

\emph{2004-2013} Associate Professor of Economic Policy, Scuola Superiore
Sant'Anna, Pisa, Italy.

\emph{since 2002} Director of the Center for the Analysis of Financial and
Economic Dynamics (CAFED), Scuola Superiore Sant'Anna Pisa, Italy.

\emph{1999-2004} Assistant Professor of Economics, Scuola Superiore
Sant'Anna, Pisa, Italy.

\emph{1998-1999} Post-doc position at the aboratory of Economic and
Management, Scuola Superiore Sant'Anna, Pisa, Italy.


\section{Education%
  \label{education}%
}

\emph{1993} Degree in Physics, University of Parma, Italy.

\emph{1998} Ph.D. in Theoretical Physics, University of Milano, Italy.


\subsection{Attended International Schools%
  \label{attended-international-schools}%
}

\emph{1998} School on the mathematics of economics, at SISSA, Trieste, Italy.

\emph{1997} Cracow School of Theoretical Physics XXXVII Course, Zakopane,
Polland.

\emph{1995} CTEQ-DESY School on Theory and Phenomenology, Bad Lauterberg,
Germany.

\emph{1995} International School of Physics 'Enrico Fermi', CXXX Course
'Selected Topics in Non Perturbative QCD', Varenna, Italy.

\emph{1994} International School of Subnuclear Physics, 32nd Course, Erice,
Italy.

\emph{1994} St.Petersburg Winter School in theoretical physics, San
Petersburg, Russia


\section{Languages%
  \label{languages}%
}

English: excellent in reading, fluent in speaking and writing

French: excellent in reading, working knowledge in speaking and basic
knowledge in writing.

Italian: native language, excellent in reading, speaking and writing


\section{Teaching Activities%
  \label{teaching-activities}%
}

\emph{1998-1999} Mathematics (Calculus) at the Department of Food Science
of the University of Parma, Italy.

\emph{1998-1999} Mathematics (Calculus and Linear Algebra) at the
Department of Social Sciences, Scuola Superiore Sant'Anna, Italy

\emph{2000-2001} Mathematics (Calculus) at the Department of Engineering of
the University of Modena and Reggio, Italy.

\emph{1999-2003} Coordinator of the courses of Mathematics and
Statistics of the PhD Program in Economics and Management of Scuola
Superiore Sant'Anna.

\emph{since 2004} Industrial Dynamics and Quantitative Methods, Graduate
Course in Economics, Scuola Superiore Sant'Anna.

\emph{2007} and \emph{2010} DIMETIC Program Lecturer, Pecs University, Hungary.

\emph{since 2008} Quantitative Methods for Economics, International PhD
Program in Economics, Scuola Superiore Sant'Anna.

\begin{description}
\item[{\emph{since 2010} Mathematics for Economics, Undegraduate and Master}] 
levels, Scuola Superiore Sant'Anna.

\end{description}

\emph{2010-2011} Microeconomics, Master Degree in Economics, University of Pisa.

\emph{since 2011} Financial Economics, Master Degree in Economics, University of Pisa.

\emph{since 2018} Financial Economics, Ph.D in Economics, Scuola Superiore Sant'Anna.


\section{Fellowships%
  \label{fellowships}%
}

Since 2003 member of the Italian Society for Chaos and Complexity.

Since 2004 member of the Italian Economic Association.

Since 2005 member of the International Schumpeterian Society.

Since 2006 member of the European Economic Association.

2007 Member of the Scientific Commission of the 'Società Italiana di
Statistica' (Italian Statistical Society) for the Economic and
Environmental Research.


\section{Scientific Activity%
  \label{scientific-activity}%
}


\subsection{Referee%
  \label{referee}%
}

Since 1998 collaboration as referee for many scientific journals,
among which:

\begin{itemize}
\item Empirical Economics

\item Europhysics Journal-B

\item International Journal of Industrial Organization

\item Industrial and Corporate Change

\item Information Sciences

\item Journal of Business and Economic Statistics

\item Journal of Economic Dynamics and Control

\item Journal of Evolutionary Economics

\item Journal of Econometrics

\item Journal of Mathematical Economics

\item Physica A

\item Research Policy

\item Review of Social Studies

\item Review of Industrial Organization

\item Regional Studies

\item Small Business Economics

\item Structural Change and Economic Dynamics
\end{itemize}


\subsection{Main event organized or advised%
  \label{main-event-organized-or-advised}%
}

September 2022, Summer School of Mathematics for Economics and Social
Sciences, Ninth Edition, San Miniato, Italy

September 2019, Summer School of Mathematics for Economics and Social
Sciences, Eight Edition, San Miniato, Italy

September 2018, Summer School of Mathematics for Economics and Social
Sciences, Seventh Edition, San Miniato, Italy

September 2017, Summer School of Mathematics for Economics and Social
Sciences, Sixth Edition, San Miniato, Italy

September 2016, Summer School of Mathematics for Economics and Social
Sciences, Fifth Edition, San Miniato, Italy

September 2015, Summer School of Mathematics for Economics and Social
Sciences, Fourth Edition, San Miniato, Italy

September 2014, Summer School of Mathematics for Economics and Social
Sciences, Third Edition, San Miniato, Italy

October 2013, \textquotedbl{}The growth of firms and countries\textquotedbl{}, PRIN final
worskshop, Rome, Italy.

September 2013, Summer School of Mathematics for Economics and Social
Sciences, Second Edition, San Miniato, Italy

December 2012, \textquotedbl{}The growth of firms and countries: distributional
properties and economic determinants\textquotedbl{}, PRIN mid-term worskshop, Pisa,
Italy.

September 2012 \textquotedbl{}Summer School of Mathematics for Economics and Social
Sciences\textquotedbl{}, First Edition, San Miniato, Italy.

October 2011, \textquotedbl{}The growth of firms and countries\textquotedbl{}, PRIN Project
kick-off meeting, Pisa, Italy.

February 2011, \textquotedbl{}EMAEE 2011: 7th European Meeting on Applied
Evolutionary Economics\textquotedbl{}, Pisa, Italy.

December 2009, \textquotedbl{}Markets and firms dynamics: from micro patterns to
aggregate behaviors\textquotedbl{}, Université des Antilles et de la Guyane,
Pointe-à-Pitre, Guadalupe, France.

October 2009 \textquotedbl{}Evolution and market behavior in economics and finance\textquotedbl{},
Mathematics Research Center \textquotedbl{}Ennio De Giorgi\textquotedbl{}, Pisa, Italy.

October 2008 \textquotedbl{}The dynamics of firm evolution: productivity,
profitability and growth\textquotedbl{}, Scuola Superiore Sant'Anna, Italy.

June 2008 \textquotedbl{}ENTREPRENEURSHIP AND INNOVATION\textquotedbl{}, 25th DRUID Celebration
Conference 2008, Copenhagen Business School, Denmark.

July 2007 \textquotedbl{}The Emergence and Impact of Market : The Case of Fish
Markets\textquotedbl{}, Tromsø, Norway.

September 2007 \textquotedbl{}Agent Based Models for Spatial Systems in Social
Sciences \& Economic Science with Heterogeneous Interacting Agents\textquotedbl{},
La Londe les Maures, France.

September 2006 \textquotedbl{}Artificial Economics 2006\textquotedbl{}, Aalborg, Denmark.

May 2006 \textquotedbl{}Complexity 2006\textquotedbl{}, Aix-en-Provence, France.

Agust 2004 IWAMEM -{}- \textquotedbl{}Data Mining and Adaptive Modeling Methods for
Economics and Management\textquotedbl{}, University of Pisa, Pisa, Italy.

April-July 2002 \textquotedbl{}Financial Markets: mathematical, statistical and
economic analysis\textquotedbl{}, Centro di Ricerca Matematica E. De Giorgi, Pisa,
Italy.


\subsection{Main attended meetings%
  \label{main-attended-meetings}%
}

June 2019 , Computing in Economics and Finance, Mila, Italy.

August 2018 invited speaker, AMMCS-2019, Waterloo, Ontario, Canada.

October 2017 speaker, Meeting of the Association of Southern-European
Economic Theorists (ASSET), Algiers, Algeria.

June 2017 invited speaker, SAET Conference, Ancao, Portugal.

June 2016 speaker, 22th Annual Workshop on Economic Heterogeneous
Interacting Agents, Universita Cattolica del Sacro Cuore , Milan.

June 2017 invited speaker, SAET Conference, Ancao, Portugal.

December 2016 invited speaker, ICC Conference: Beyond Technological
Innovation and Diffusion, University of California, Berkeley, US.

October 2016 speaker, Meeting of the Association of Southern-European
Economic Theorists (ASSET), Thessaloniki, Greece.

June 2016 speaker and invited teacher, 21th Annual Workshop on
Economic Heterogeneous Interacting Agents, University of Castellon,
Spain.

October 2015 speaker, Meeting of the Association of Southern-European
Economic Theorists (ASSET), Grenada, Spain.

July 2015 invited speaker, IWH Workshop on Firm Exit and Job
Displacement, Halle Institute for Economic Research, Halle, Germany.

May 2015 speaker, 20th Annual Workshop on Economic Heterogeneous
Interacting Agents, University of Nice, France.

June 2014 speaker, Governance of a Complex World 2014, Torino, Italy.

May 2013 speaker and invited teacher, 18th Annual Workshop on
Economic Heterogeneous Interacting Agents, University of Reykjavik,
Island.

June 2012 speaker and tutorial teacher, 17th Annual Workshop on
Economic Heterogeneous Interacting Agents, University of
Pantheon-Assas, Paris II, France.

February 2012 invited speaker, Financing Innovation and Growth:
Reforming a Dysfunctional System, London, UK

February 2012 invited speaker, External Developments Division of the
European Central Bank, Frankfurt, Germany.

June 2011 invited speaker, SAET Conference, Ancao, Portugal.

June 2011 invited speaker, International Conference STOCHASTIC
ECONOMICS AND FINANCE, University of Bergen, Bergen, Norway.

January 2011 invited speaker, Second International Meeting of TOM
Society, Universitat Pompeu Fabra, Barcelona, Spain.

October 2010 invited speaker \textquotedbl{}Quantifying and Understanding
Dysfunctions of Financial Markets\textquotedbl{}, Luvein, Belgium

September 2010 invited speaker \textquotedbl{}Toward an Alternative Macroeconomic
Theory\textquotedbl{} Central European University, Budapest

June 2010 \textquotedbl{}11th Workshop on Optimal Control, Dynamic Games and
Nonlinear Dynamics\textquotedbl{}, University of Amsterdam, The Netherlands.

June 2010 \textquotedbl{}International Schumpeter Conference\textquotedbl{}, Aalborg University,
Denmark.

June 2009 invited speaker \textquotedbl{}Common Complex Collective Phenomena:
implications for economic and social policy making\textquotedbl{}, Cracow, Polland.

January 2009 invited speaker at the Workshop \textquotedbl{}Ten Years of CeNDEF\textquotedbl{},
University of Amsterdam, The Netherlands.

November 2008 invited speaker at the Workshop \textquotedbl{}Revolving Doors:
Entrepreneurial Survival and Exit\textquotedbl{}, Rotterdam School of Management,
The Netherlands.

October 2008 invited speaker at the 3rd Toyota CRDL Workshop
\textquotedbl{}Mathematical Methods in Complex Systems\textquotedbl{}, Gemenos, France.

August 2008 EEA-ESEM 2008, Milano, Italy.

June 2008 \textquotedbl{}Computing in Economics and Finance\textquotedbl{}, Paris, France.

April 2008 session organizer \textquotedbl{}DIME Scientific Conference\textquotedbl{}, BETA/PEGE,
Strasbourg, France.

February 2008 invited speaker \textquotedbl{}FIG workshop\textquotedbl{}, CIRCLE, Lund, Sweden.

November 2007 invited speaker 'Second European PhD Complexity School',
Torino, Italy.

October 2007 invited speaker 'S.I.E. Annual Conference', Torino,
Italy.

October 2007 invited speaker \textquotedbl{}Complexity in Economics and Finance\textquotedbl{},
Leiden University, the Netherland.

August 2007  EEA-ESEM 2007, Budapest, Hungary.

September 2006 invited speaker \textquotedbl{}Policy Implications from Recent
Advances in the Economics of Innovation and Industrial Dynamics\textquotedbl{}, Open
University, London.

June 2006 'International Schumpeter Conference', Nice, France.

May 2006 'Complexity 2006', Aix-en-Provence, France.

September 2005 'Artificial Economics 2005', Lille, France.

May 2005 invited speaker 'Third SPIE International Symposium
Fluctuations and Noise' Austin, Texas

November 2004 invited speaker 'Workshop On Interactions and
Markets', Pisa, Italy.

June 2004 invited speaker 'Workshop On Nonlinear Dynamics in
Economics' Trieste, Italy.

June 2004 'Schumpeter 2004 Conference', Milano, Italy.

October 2003 invited speaker 'Workshop on Economic Dynamics',
Bolzano, Italy.

October 2003 \href{mailto:'Wild@Ace}{'Wild@Ace}' Torino, Italy.

September 2003 '13th European Colloquium on Quantitative and
Theoretical Geography', Lucca, Italy.

September 2003 'Clusters, Industrial Districts and Firms. The
Challenge of Globalization' Modena, Italy.

April 2003 'European Meeting on Applied Evolutionary Economics',
Augsburg, Germany.

April 2003 'Reinventing Regions in the Global Economy', Pisa, Italy.

March 2003 invited speaker 'The Evolution of Innovation and Market
Structure in the Pharmaceutical and Biotech Industry: Alternative
Models and Interpretations', Londra, United Kingdom.

March 2003 '11th Annual Symposium of the Society for Non Linear
Dynamics and Econometrics' SNDE-2003, Firenze, Italy.

September 2002 EARIE'29 29-th Annual Conference of the European
Association for Research in Industrial Economics' Madrid, Spain.

June 2002 'Evolutionary Finance Conference', Zurich, Switzerland.

May 2002 WHEIA2002 'Workshop on Economics with Heterogeneous
Interacting Agents', Trieste, Italy.

April 2002 'Local Models of Development, Clusters of Firms and Local
Innovative Systems' Varese, Italy.

September 2001 'European Meeting on Applied Evolutionary Economics',
Vienna, Austria.

June 2001 'Society for Computational Economics 2001 Conference' Yale
University, New Haven, Connecticut.

June 2001 'Complexity and Industrial Clusters' , Milan, Italy.

June 2001 WHEIA2001 'Workshop on Economics with Heterogeneous
Interacting Agents', Maastricht, Netherland.

June 2000 'Third ESSY workshop', Berlin, Germany.

June 1999 WHEIA1999 'Workshop on Economics with Heterogeneous
Interacting Agents', Genova, Italy.

Luglio 1998, 'ICHEP 98, XXIX International Conference on High Energy
Physics' Vancouver, Canada.

Giugno 1998, 'Low x Physics at HERA', DESY Zheuten, Berlin, Germany.

Settembre 1996, V National Seminary of Theoretical Physics, Parma,
Italy.

Maggio 1996, 'Cortona 96: Convegno Nazionale Informale di Fisica
Teorica ' ( Cortona 96: National Informal Conference of Theoretical
Physics), Cortona.


\subsection{Grants received%
  \label{grants-received}%
}

OPENMUSE

\begin{quote}
\begin{DUfieldlist}
\item[{period:}]
2022-2024

\item[{type:}]
Horizon 2022

\item[{sponsor:}]
European Commission

\item[{role:}]
Research group coordinator
\end{DUfieldlist}
\end{quote}

DOLFINS

\begin{quote}
\begin{DUfieldlist}
\item[{period:}]
2015-2018

\item[{type:}]
Research and Innovation Action project FP 7

\item[{sponsor:}]
European Commission

\item[{role:}]
Member
\end{DUfieldlist}
\end{quote}

Economia dei Media Sociali: Informazione e Reputazione nella Rete

\begin{quote}
\begin{DUfieldlist}
\item[{period:}]
2013-2015

\item[{type:}]
FORTEC Research Grant

\item[{sponsor:}]
Tuscany Regional Council

\item[{role:}]
Lead Scientist
\end{DUfieldlist}
\end{quote}

Market Selection and Aggregate Economic Outcomes

\begin{quote}
\begin{DUfieldlist}
\item[{period:}]
2013-2015

\item[{type:}]
MARIE CURIE INTERNATIONAL OUTGOING FELLOWSHIPS

\item[{sponsor:}]
European Commission

\item[{role:}]
Primary Coordinator
\end{DUfieldlist}
\end{quote}

The growth of firms and countries

\begin{quote}
\begin{DUfieldlist}
\item[{period:}]
2011-2013

\item[{type:}]
PRIN (Project of Relevant National Interest)

\item[{sponsor:}]
Italian Ministry of University and Research

\item[{role:}]
National Coordinator
\end{DUfieldlist}
\end{quote}

Micro-dyn

\begin{quote}
\begin{DUfieldlist}
\item[{period:}]
2006-2010

\item[{type:}]
Integrated Project in FP 6

\item[{sponsor:}]
European Commission

\item[{role:}]
External Expert
\end{DUfieldlist}
\end{quote}

Dissemination to and Interaction with Industry and other stakeholders

\begin{quote}
\begin{DUfieldlist}
\item[{period:}]
2007-2009

\item[{pype:}]
Work Package, DIME Network of Excellence

\item[{sponsor:}]
European Commission

\item[{role:}]
Coordinator
\end{DUfieldlist}
\end{quote}

CO3 - Common Complex Collective Phenomena in Statistical Mechanics ,
Society, Economics and Biology

\begin{quote}
\begin{DUfieldlist}
\item[{period:}]
2005-2009

\item[{type:}]
Specific Targeted Research Project

\item[{sponsor:}]
European Commission

\item[{role:}]
Member
\end{DUfieldlist}
\end{quote}

Market Dynamics: analytical tools, empirical regularities and
experimental results

\begin{quote}
\begin{DUfieldlist}
\item[{period:}]
2006-2008

\item[{type:}]
Work Package, DIME Network of Excellence

\item[{sponsor:}]
European Commission

\item[{role:}]
Coordinator
\end{DUfieldlist}
\end{quote}

Apprendimento e interazioni compettive nei modelli di evoluzione delle
insustrie

\begin{quote}
\begin{DUfieldlist}
\item[{period:}]
2004-2006

\item[{type:}]
PRIN

\item[{sponsor:}]
Italian Ministry of University and Research

\item[{role:}]
Member
\end{DUfieldlist}
\end{quote}

Modelli dinamici dei mercati finanziari e analisi delle serie storiche
ad alta frequenza: dai processi stocastici alla gestione del rischio

\begin{quote}
\begin{DUfieldlist}
\item[{period:}]
2002-2003

\item[{type:}]
Programma di Ricerca di Rilevante Interesse

\item[{sponsor:}]
Scuola Superiore Sant'Anna

\item[{role:}]
Coordinator
\end{DUfieldlist}
\end{quote}

Mercati ed istituzioni: dinamiche di apprendimento e strutture di
interazione

\begin{quote}
\begin{DUfieldlist}
\item[{period:}]
2001-2002

\item[{type:}]
Programma di Ricerca di Rilevante Interesse

\item[{sponsor:}]
Scuola Superiore Sant'Anna

\item[{role:}]
Member
\end{DUfieldlist}
\end{quote}

Structures and dynamics of industries and markets

\begin{quote}
\begin{DUfieldlist}
\item[{period:}]
2001-2002

\item[{type:}]
FIRB

\item[{sponsor:}]
Italian Ministry of University and Research

\item[{role:}]
Member
\end{DUfieldlist}
\end{quote}

Metodi e modelli statistici per dati macroeconomici e industriali:
analisi sezionali e temporali

\begin{quote}
\begin{DUfieldlist}
\item[{period:}]
2002

\item[{type:}]
Coordinated project - AGENZIA 2000

\item[{sponsor:}]
CNR

\item[{role:}]
Member
\end{DUfieldlist}
\end{quote}

Modelli dinamici per la crescita e diversificazione dell'impresa

\begin{quote}
\begin{DUfieldlist}
\item[{period:}]
1999-2001

\item[{type:}]
Progetti Giovani Ricercatori

\item[{sponsor:}]
CNR

\item[{role:}]
Coordinator
\end{DUfieldlist}
\end{quote}


\section{Main publications in Economics%
  \label{main-publications-in-economics}%
}


\subsection{Publications in international journals%
  \label{publications-in-international-journals}%
}

G.Bottazzi and D. Giachini, A general equilibrium model of investor
sentiment, Economics Letters, Volume 218, September 2022.

G. Bottazzi and D. Giachini, STRATEGICALLY BIASED LEARNING IN MARKET
INTERACTIONS, Advances in Complex Systems, Vol. 25, No. 02n03,
2250004, 2022.

G. Bottazzi, On the Pareto Type III Distribution, Frontiers in
Physics, Sec. Interdisciplinary Physics, 14 April 2022.

G. Bottazzi and P. Dindo, Drift criteria for persistence of discrete
stochastic processes on the line, Journal of Mathematical Economics,
Volume 101, 2022.

G. Bottazzi, Dindo P., Giachini D., Momentum and reversal in financial
markets with persistent heterogeneity, Annals of Finance, vol. 15(4),
pages 455-487, 2019.

G. Bottazzi, Li L., Secchi A. . Aggregate fluctuations and the
distribution of firm growth rates, Industrial and Corporate Change,
vol. 28, p. 635-656, 2019

G. Bottazzi and D. Giachini, Far from the Madding Crowd: Collective
Wisdom in Prediction Markets, Quantitative Finance,
vol. 19, p. 1461-1471 , 2019.

G. Bottazzi and D. Giachini, Betting, Selection, and Luck: A Long-Run
Analysis of Repeated Betting Markets, Entropy 21(6), 585, June, 2019.

G. Bottazzi, P. Dindo and D. Giachini, Long-run Heterogeneity in an
Exchange Economy with Fixed-Mix Traders, Economic Theory, Volume 66,
pp 407-447, August 2018.

G. Bottazzi and D. Giachini Wealth and price distribution by diffusive
approximation in a repeated prediction market, Physica A, Volume 471,
pp 473-479, 1 April 2017.

G. Bottazzi, U. Gragnolati and F. Vanni Non-linear externalities in
firm localization Regional Studies, November 2016.

S. Bianchini, G. Bottazzi and F. Tamagni, What does (not) determine
persistent corporate high-growth? Small Business Economics,
July, 2016.

G. Bottazzi, D. Pirino and F. Tamagni, Zipf law and the firm size
distribution: a critical discussion of popular estimators, Journal of
Evolutionary Economics, 25 (3), 585-610, 2015.

G. Bottazzi and U. Gragnolati, Cities and clusters: economy-wide and
sector-specific effects in corporate location, Regional Studies 49
(1),pp 113-129, 2015.

G. Bottazzi and M. Grazzi, Dynamics of productivity and cost of labor:
Evidence from Italian manufacturing firms, Bulletin of Economic
Research, DOI:10.1111/boer.12014

G. Bottazzi and P. Dindo, Globalizing knowledge: how technological
openness affects output, spatial inequality and welfare levels,
Journal of Regional Science, DOI:10.1111/jors.12034

G. Bottazzi and U. Gragnolati, Cities and clusters: economy-wide and
sector-specific effects in corporate location, Regional Studies,
DOI:10.1080/00343404.2012.739281

G. Bottazzi, A. Secchi, F. Tamagni, Financial Constraints and Firm
Dynamics, Small Business Economics, DOI:10.1007/s11187-012-9465-5

G. Bottazzi and P. Dindo, Selection in asset markets: the good, the
bad, and the unknown, Journal of Evolutionary Economics, Volume 23,
Issue 3, pp 641-661, July 2013.

G. Bottazzi and P. Dindo, Evolution and Market Behavior in Economics
and Finance: introduction to the special issue, Journal of
Evolutionary Economics, Volume 23, Issue 3, pp 507-512, July 2013

M. Anufriev and G. Bottazzi, Asset Pricing with Heterogeneous
Investment Horizons, Studies in Nonlinear Dynamics \& Econometrics,
16(4), 2012.

G. Bottazzi and A. Secchi, Productivity, profitability and growth: the
empirics of firm dynamics, Structural Change and Economic Dynamics,
23(4), pp.325–328, 2012.

M. Anufriev, G. Bottazzi, M. Marsili and P. Pin, Excess Covariance and
Dynamic Instability in a Multi-Asset Model, Journal of Economic
Dynamics and Control, 36(8), pp.1142-1161, 2012.

G. Bottazzi and F. Tamagni, Big and fragile: when size does not shield
from defaults, Applied Economics Letters, 18(14), pp. 1401-1404, 2011.

G. Bottazzi and A. Secchi, A New Class of Asymmetric Exponential Power
Densities with Applications to Economics and Finance, Industrial and
Corporate Change, 20(4), pp. 991-1030, 2011.

G. Bottazzi, M. Grazzi, A. Secchi, F. Tamagni, Financial and Economic
Determinants of Firm Default, Journal of Evolutionary Economics,
21(3), pp. 373-406, 2011.

G. Bottazzi, G. Devetag, F. Pancotto, Does Volatility matter?
Expectations of price return and variability in an asset pricing
experiment, Journal of Economic Behavior \& Organization, 77,
pp. 124-146, 2011.

G.Bottazzi, A.Coad, N.Jacoby and A.Secchi, Corporate growth and
industrial dynamics: evidence from French manufacturing, Applied
Economics, 43 (1), pp.103-116, 2011.

M. Anufriev and G. Bottazzi, Market Equilibria under Procedural
Rationality, Journal of Mathematical Economics, 46 (6), pp. 1140–1172,
2010.

G. Bottazzi, G. Dosi, N. Jacoby, A. Secchi, F. Tamagni, Corporate
performances and market selection. Some comparative evidence,
Industrial and Corporate Change, 19 (6), pp. 1953-1996, 2010.

G.Bottazzi and M.Grazzi, Wage-size relation and the structure of
work-force composition in Italian Manufacturing Firms, Cambridge
Journal of Economics, Volume 34, Number 4 pp. 649-669, 2010.

G.Bottazzi, On the Irreconcilability of Pareto and Gibrat Laws,
Physica A, 388, 7, pp. 1133-1136, 2009.

G.Bottazzi, A.Secchi and F.Tamagni, Productivity, Profitability and
Financial Performance, Industrial and Corporate Change, 17,
pp.711-751, 2008.

G.Bottazzi, G.Fagiolo, G.Dosi and A.Secchi, Sectoral and Geographical
Specificities in the Spatial Structure of Economic Activities,
Structural Change and Economic Dynamics, 19, pp. 189-202, 2008.

G.Bottazzi, On the relationship between firms' size and growth rate,
Economics Bulletin, 3, 8, pp. 1-7, 2008.

G.Bottazzi, G. Fagiolo, G.Dosi and A. Secchi, Modeling Industrial
Evolution in Geographical Space, Journal of Economic Geography,
7, pp. 651-672, 2007.

G.Bottazzi and G.Devetag, Competition and Coordination in Experimental
Minority Games, Journal of Evolutionary Economics, 17, pp. 241-275,
2007.

G.Bottazzi, E.Cefis, G.Dosi and A.Secchi, Invariances and Diversities
in the Evolution of Italian Manufacturing Industry, Small Business
Economics, 29, pp. 137-159, 2007.

G.Bottazzi and A.Secchi, Gibrat's Law and Diversification, Industrial
and Corporate Change, 15, pp. 847-875, 2006.

M.Anufriev, G.Bottazzi and F.Pancotto, Equilibria, Stability and
Asymptotic Dominance in a Speculative Market with Heterogeneous
Agents, Journal of Economic Dynamics and Control, 30, pp. 1787-1835,
2006.

G.Bottazzi and A.Secchi, Explaining the Distribution of Firms Growth
Rates, Rand Journal of Economics, vol. 37, no. 3, 2006.

G.Bottazzi, M.Grazzi and A.Secchi, Characterizing the Production
Process: A Disaggregated Analysis of Italian Manufacturing Firms,
Rivista di Politica Economica, I-II,pp.  243-270, 2005.

G.Bottazzi and A.Secchi, Growth and Diversification Patterns of the
Worldwide Pharmaceutical Industry, Review of Industrial Organization,
26, pp. 195-216, 2005.

G.Bottazzi, G.Dosi and I.Rebesco, Institutional Architectures and
Behavioral Ecologies in the Dynamics of Financial Markets, Journal of
Mathematical Economics, vol. 41, pp.  197-228, 2005.

G.Bottazzi, M.Grazzi and A.Secchi, Input Output Scaling Relations in
Italian Manufacturing Firms, Physica A, 355, pp.  95-102, 2005.

G.Bottazzi, S.Sapio and A.Secchi, Some Statistical Investigations on
the Nature and Dynamics of Electricity Prices, Physica A, 355, pp.
54-61, 2005.

G.Bottazzi, A.Secchi, Common Properties and Sectoral Specificities in
the Dynamics of U.S. Manufacturing Companies, Review of Industrial
Organization, vol. 23, pp. 217-232, 2003.

G.Bottazzi, A.Secchi, Why are distributions of firm growth rates
tent-shaped , Economic Letters vol. 80 pp.415-420, 2003.

G.Bottazzi, G. Devetag, A Laboratory Experiment on the Minority Game,
Physica A vol. 324 pp. 124-132, 2003.

G.Bottazzi, A. Secchi, A Stochastic Model of Firm Growth, Physica A
vol. 324 pp. 213-219, 2003.

G.Bottazzi, G.Devetag, G.Dosi, Learning and emergent coordination in
speculative markets: some properties of \textquotedbl{}minority game\textquotedbl{} dynamics,
Simulation Modeling Practice and Theory, 10, 321-347 (2002).

G.Bottazzi, E.Cefis, G.Dosi, Corporate Growth and Industrial
Structure. Some Evidence from the Italian Manufacturing Industry,
Industrial and Corporate Change 2,4 (2002).

G.Bottazzi, G.Dosi and G.Rocchetti, Modes of Knowledge Accumulation,
Entry Regimes and Patterns of Industrial Evolution, Industrial and
Corporate Change 10,3 (2001)

G.Bottazzi, G.Dosi, M.Lippi, F.Pammolli, M.Riccaboni, Innovation and
Corporate Growth in the Evolution of the Drug Industry, International
Journal of Industrial Organization 19 (2001) pp.1161-1187


\subsection{Articles in Monograph and Collections%
  \label{articles-in-monograph-and-collections}%
}

G.Bottazzi and A.Secchi, Financial fragility and the distribution of
firm growth rates, in A.M.Ferragina and E. Taymaz (eds) \textquotedbl{}Innovation,
Globalization and Firm Dynamics: Lessons for Enterprise Policy\textquotedbl{}, 2013.

G.Bottazzi and P.Dindo, An evolutionary model of firms location with
technological externality, in R.A.Boschma and R.Martin (eds.)
\textquotedbl{}Handbook on Evolutionary Economic Geography\textquotedbl{}, Edward Elgar,
Cheltenham, 2010.

G.Bottazzi and M.Grazzi, Dinamiche della produttivit`a e del costo
del lavoro nelle imprese manifatturiere italiane, in Produttivit`a e
cambiamento nell'industria italiana, Laura Rondi e Francesco Silva
(eds.), il Mulino, 2009.

G.Bottazzi and S.Sapio, Power Exponential Price Returns in Day-ahead
Power Exchanges, in A. Chatterjee and B. K. Chakrabarti (eds)
\textquotedbl{}Econophysics of Markets and Business Networks\textquotedbl{}, New Economic Windows,
Springer Verlag, Berlin, 2007.

G.Bottazzi, M. Grazzi and A. Secchi, Characterizing the Production
Process: A Disaggregated Analysis of Italian Manufacturing Firms, in
M. Malgarini and G. Piga (eds) \textquotedbl{}Capital Accumulation, Productivity and
Growth\textquotedbl{}, Palgrave Macmillan, New York, 2006.

G.Bottazzi and A. Secchi, Self-reinforcing Dynamics and the evolution
of business firms, in A.Pyka and and H.Hanusch (eds.) \textquotedbl{}Applied
Evolutionary Economics and the Knowledge-Based Economy\textquotedbl{}, Edward Elgar,
Cheltenham, 2006.

M.Anufriev and G. Bottazzi, Noisy Trading in the Large Market Limit,
in P.Mathieu, B.Beaufils and O.Brandouy (eds.)  \textquotedbl{}Artificial Economics:
Agent-based Methods in Finance, Game Theory and Their Applications\textquotedbl{},
Lecture Notes in Economics and Mathematical Systems vol.564,
Springer-Verlag, Berlin, 2006.

G.Bottazzi, G Dosi, and G. Fagiolo, On Sectoral Specificities in the
Geography of Corporate Location, in Breschi, S. and Malerba, F.
(Eds.) \textquotedbl{}Clusters, networks and innovation\textquotedbl{}, Oxford University Press,
Oxford, U.K., 2006.

G.Bottazzi, F.Pammolli and A.Secchi, The Dynamics of Growth and
Diversification of the Large Pharmaceutical Companie,, in M.
Mazzucato and G.  Dosi (Eds.) 'Knowledge Accumulation and Industry
Evolution: the Case of Pharma-Biotech', Cambridge University Press,
2005.

G.Bottazzi, M. G. Devetag, Expectations structure in asset pricing
experiments, in T.Lux, S.Reitz and E.Samanidou (Eds.)  \textquotedbl{}Nonlinear
Dynamics and Heterogeneous Interacting Agents\textquotedbl{} Lecture Notes in
Economics and Mathematical Systems 550, Springer Verlag, Berlin, 2005.

G.Bottazzi, M. G. Devetag, Coordination and Self-Organization in
Minority Games: Experimental Evidence, in M.Gallegati A. Kirman and
M.Marsili (Eds.) \textquotedbl{}The Complex Dynamics of Economic Interaction\textquotedbl{},
Lecture Notes in Economics and Mathematical Systems (531) pp. 283-300,
Springer Verlag, Berlin, 2004.

G.Bottazzi, A Simple Micro-Model of Market Dynamics, in
M.Gallegati A. Kirman and M.Marsili (Eds.) \textquotedbl{}The Complex Dynamics of
Economic Interaction\textquotedbl{}, Lecture Notes in Economics and Mathematical
Systems (531) pp. 283-300, Springer Verlag, Berlin, 2004.

G.Bottazzi, G.Devetag, G.Dosi, Adaptive Learning and Emergent
Coordination in Minority Games, in R.Cowan and N.Jonard Eds.
\textquotedbl{}Heterogeneous Agents, Interactions and Economic Performance\textquotedbl{} Lecture
Notes in Economics and Mathematical Systems 521, Springer Verlag,
Berlin, 2003.

G.Bottazzi, E.Cefis, G.Dosi e A.Secchi, Crescita dell'impresa e
struttura industriale: evidenze empiriche sull'industria
manifatturiera italiana, in D. Delli Gatti e M.Gallegati (eds.)
Eterogeneita' degli agenti economici ed interazione sociale: teorie e
verifiche empiriche, Il Mulino, Bologna, 2003.

G.Bottazzi, G.Fagiolo, G.Dosi, On the ubiquitous nature of the
agglomeration economies and their limits, in A. Quadro Curzio and
M.Fortis eds. 'Complexity and Industrial Cluster', Physica-Verlag,
Heidelberg, 2002.

G.Bottazzi, G.Fagiolo, G.Dosi, Sulla onnipresenza delle economie di
agglomerazione e le loro diverse determinanti: alcune note, in A.
Quadro Curzio and M. Fortis (a cura di) 'Complessita' e distretti
industriali: dinamicge, modelli, casi reali', Il Mulino Collana della
Fondazione Edison, pp. 265-298, Bologna, 2002.


\subsection{Referred Contributions to Proceedings%
  \label{referred-contributions-to-proceedings}%
}

M.Anufriev, G. Bottazzi and F. Pancotto, Speculative Equilibria and
Asymptotic Dominance in a Market with Adaptive CRRA Traders, in Abbot,
D., Bouchard, J.P., Gabaix, X., and McCauley, J.L. \textquotedbl{}Noise and
Fluctuations in Econophysics and Finance\textquotedbl{} SPIE Conference Proceedings,
vol. 5848, SPIE, WA, 2005.


\subsection{Editorial work%
  \label{editorial-work}%
}

Evolution and market behavior in economics and finance, Giulio
Bottazzi and Pietro Dindo (eds), Journal of Evolutionary Economics,
Volume 23(3), 2013.

Productivity, profitability and growth: the empiric of firm dynamics,
Giulio Bottazzi and Angelo Secchi (eds.), Industrial and Corporate
Change 23(4), 2012.


\subsection{Working Papers%
  \label{working-papers}%
}

G. Bottazzi, D. Giachini, New Results on Betting Strategies, Market
Selection, and the Role of Luck, L.E.M.  Working Paper n. 2016-14,
Scuola Superiore Sant'Anna, (2018).

G. Bottazzi, P. Dindo, D. Giachini, Momentum and Reversal in Financial
Markets with Persistent Heterogeneity, L.E.M.  Working
Paper n. 2016-14, Scuola Superiore Sant'Anna, (2018).

G. Bottazzi, L. Li, A. Secchi, Aggregate fluctuations and the
distribution of firm growth rates, L.E.M.  Working Paper n. 2016-14,
Scuola Superiore Sant'Anna, (2017).

G. Bottazzi, D. Giachini, Far from the Madding Crowd: Collective
Wisdom in Prediction Markets, L.E.M.  Working Paper n. 2016-14, Scuola
Superiore Sant'Anna, (2016).

G. Bottazzi, D. Giachini, Wealth and Price Distribution by Diffusive
Approximation in a Repeated Prediction Market, L.E.M.  Working
Paper n. 2016-13, Scuola Superiore Sant'Anna, (2016).

G.Bottazzi, A.De Sanctis, F.Vanni, Non-performing loans, systemic risk
and resilience in financial networks, L.E.M.  Working
Paper n. 2016-08, Scuola Superiore Sant'Anna, (2016).

G.Bottazzi, P.Dindo, D.Giachini, Long-run Heterogeneity in an Exchange
Economy with Fixed-Mix Traders, L.E.M.  Working Paper n. 2015-29,
Scuola Superiore Sant'Anna, (2015).

G.Bottazzi, U.Gragnolati, F.Vanni, Non-linear externalities in firm
localization, L.E.M.  Working Paper n. 2015-28, Scuola Superiore
Sant'Anna, (2015).

G.Bottazzi, P.Dindo, Drift criteria for persistence of discrete
stochastic processes on the line, L.E.M.  Working Paper n. 2015-26,
Scuola Superiore Sant'Anna, (2015).

S. Bianchini, G. Bottazzi, F. Tamagni, What does (not) determine
persistent corporate high-growth? L.E.M.  Working Paper n. 2014-11,
Scuola Superiore Sant'Anna, (2014).

G. Bottazzi, F. Vanni, A numerical estimation method for discrete
choice models with non-linear externalities L.E.M.  Working
Paper n. 2014-1, Scuola Superiore Sant'Anna, (2014).

G.Bottazzi, D.Pirino, F.Tamagni, Zipf Law and the Firm Size
Distribution: a critical discussion of popular estimators, L.E.M.
Working Paper n. 2013-17, Scuola Superiore Sant'Anna, (2013).

G. Bottazzi, A. Secchi, Financial Fragility and the Distribution of
Firm Growth Rates, L.E.M.  Working Paper n. 2013-13, Scuola Superiore
Sant'Anna, (2013).

G. Bottazzi, M. Grazzi, Dynamics of productivity and cost of labor in
Italian Manufacturing firms, L.E.M.  Working Paper n. 2013-02, Scuola
Superiore Sant'Anna, (2013).

G. Bottazzi, M. Duenas, The Evolution of the Business Cycles and
Growth Rates Distributions, L.E.M.  Working Paper n. 2012-22, Scuola
Superiore Sant'Anna, (2012).

G. Bottazzi, A.Secchi, A New Class of Asymmetric Exponential Power
Densities with Applications to Economics and Finance, L.E.M.  Working
Paper n. 2012-22, Scuola Superiore Sant'Anna, (2011).

M. Anufriev, G. Bottazzi, M. Marsili, P. Pin, Excess Covariance and
Dynamic Instability in a Multi-Asset Model, CeNDEF Working Papers
11-09, (2011)

G. Bottazzi P. Dindo Selection in asset markets: the good, the bad,
and the unknown, L.E.M.  Working Paper n. 2010-20, Scuola Superiore
Sant'Anna, (2011).

G. Bottazzi, P. Dindo, Evolution and market behavior with endogenous
investment rules, L.E.M.  Working Paper n. 2010-20, Scuola Superiore
Sant'Anna, (2010).

G. Bottazzi, U. Gragnolati, Cities and clusters: economy-wide and
sector-specific effects in corporate location, L.E.M.  Working Paper
n. 2010-16, Scuola Superiore Sant'Anna, (2010).

G. Bottazzi, D. Pirino, Measuring Industry Relatedness and Corporate
Coherence, L.E.M.  Working Paper n. 2010-10, Scuola Superiore
Sant'Anna, (2010).

G. Bottazzi, F. Tamagni, Is Bigger Always Better ? The effect of Size
on Defaults, L.E.M.  Working Paper n. 2010-07, Scuola Superiore
Sant'Anna, (2010).

G. Bottazzi, M. Grazzi, A. Secchi, F. Tamagni, Financial and Economic
Determinants of Firm Default, L.E.M.  Working Paper n. 2009-06, Scuola
Superiore Sant'Anna, (2009).

G.Bottazzi, G.Devetag and F.Pancotto, Does Volatility matter?
Expectations of price return and variability in an asset pricing
experiment,L.E.M.  Working Paper n. 2009-02, Scuola Superiore Sant'Anna,
(2009).

G. Bottazzi and P. Dindo, An evolutionary model of firms location with
technological externalities, L.E.M.  Working Paper n. 2008-27, Scuola
Superiore Sant'Anna (2008).

G.Bottazzi and P.Dindo, Localized technological externalities and the
geographical distribution of firms, L.E.M.  Working Paper n. 2008-11,
Scuola Superiore Sant'Anna (2008).

G.Bottazzi, M.Grazzi, A.Secchi, F.Tamagni, Assessing the Impact of
Credit Ratings and Economic Performance on Firm Default,
L.E.M. Working Paper n. 2007-15, Scuola Superiore Sant'Anna (2007).

G.Bottazzi, A comment on the relationship between firms' size and
growth rate, L.E.M. Working Paper n. 2007-11, Scuola Superiore
Sant'Anna (2007).

G.Bottazzi, On the Irreconcilability of Pareto and Gibrat Laws,
L.E.M. Working Paper n. 2007-10, Scuola Superiore Sant'Anna (2007).

G.Bottazzi, On the Pareto Type III Distribution, L.E.M. Working Paper
n. 2007-08, Scuola Superiore Sant'Anna (2007).

G.Bottazzi, G. Dosi, G. Fagiolo and A. Secchi, Modeling Industrial
Evolution in Geographical Space, L.E.M. Working Paper n. 2007-06,
S.Anna School (2007).

G.Bottazzi and M. Grazzi, Wage structure in Italian Manufacturing
firms, L.E.M. Working Paper n. 2007-05, Scuola Superiore Sant'Anna
(2007).

G.Bottazzi, M.Grazzi and A.Secchi, Productivity, Profitability and
Financial Fragility: Empirical Evidence from Italian Business Firms,
L.E.M. Working Paper n.2006-08, Scuola Superiore Sant'Anna (2006).

G.Bottazzi, M.Grazzi and A.Secchi, Financial Fragility and Growth
Dynamics of Italian Business Firms, L.E.M. Working Paper n.2006-07,
S.Anna School (2006).

G.Bottazzi, A. Coad, N. Jacoby and A. Secchi, Corporate Growth and
Industrial Dynamics: Evidence from French Manufacturing,
L.E.M. Working Paper n.2005-21, Scuola Superiore Sant'Anna (2005).

G.Bottazzi and A. Secchi, Explaining the Distribution of Firms Growth
Rates, L.E.M. Working Paper n. 2005-16, Scuola Sant'Anna, Pisa (2005).

M.Anufriev and G.Bottazzi, Price and Wealth Dynamics in a Speculative
Market with an Arbitrary Number of Generic Technical Traders,
L.E.M. Working Paper n. 2005-16, Scuola Superiore Sant'Anna, Pisa
(2005).

G.Bottazzi, M. Grazzi and A. Secchi, Characterising the Production
Process: A Disaggregated Analysis of Italian Manufacturing Firms,
L.E.M. Working Paper n. 2004-24, Scuola Sant'Anna, Pisa.

G.Bottazzi, M. Anoufriev and F. Pancotto, Price and Wealth Asymptotic
Dynamics with CRRA Technical Trading Strategies, L.E.M. Working Paper
n. 2004-23, Scuola Superiore Sant'Anna, Pisa.

G.Bottazzi and M. Anoufriev, Asset Pricing Model with Heterogeneous
Investment Horizons, L.E.M. Working Paper n. 2004-22, Scuola Superiore
Sant'Anna, Pisa.

G.Bottazzi, G. Dosi, G. Fagiolo, A. Secchi, Sectoral and Geographical
Specificities in the Spatial Structure of Economic Activities,
L.E.M. Working Paper n. 2004-21, Scuola Superiore Sant'Anna, Pisa.

G.Bottazzi, Subbotools User's Manual, L.E.M. Working Paper n.
2004-14, Scuola Superiore Sant'Anna, Pisa.

G.Bottazzi, A.Sapio and A.Secchi, Some Statistical Investigations on
the Nature and Dynamics of Electricity Prices,, L.E.M. Working Paper
n. 2004-13, Scuola Superiore Sant'Anna, Pisa.

G.Bottazzi and A.Secchi, On the Laplace Shape of the
Distribution of Firms Growth Rates, Working Paper n. 41, School of
Economics - Free University of Bozen, Bolzano (2004).

G.Bottazzi, E.Cefis, G.Dosi and A.Secchi, Invariances and
Diversities in the Evolution of Manufacturing Industries, L.E.M.
Working Paper n. 2003-21, Scuola Superiore Sant'Anna (2003).

G.Bottazzi and G. Devetag, Expectations structure in asset
pricing experiments, L.E.M. Working Paper n. 2003-19, S.Anna
School (2003).

G.Bottazzi and A.Secchi, Sectoral Specifities in the
Dynamics of U.S. Manufacturing Firms, L.E.M. Working Paper n.
2003-18, Scuola Superiore Sant'Anna (2003).

G.Bottazzi, G.Dosi and I.Rebesco, Institutional Architectures
and Behavioral Ecologies in the Dynamics of Financial Markets: a
Preliminary Investigation, L.E.M. Working Paper n. 2002-24, S.Anna
School (2002).

G.Bottazzi, G.Dosi and G.Fagiolo, Mapping Sectoral Patterns
of Technological Accumulation into the Geography of Corporate
Locations. A Simple Model and Some Promising Evidence, L.E.M.
Working Paper n. 2002-21, Scuola Superiore Sant'Anna (2002).

G.Bottazzi, A. Secchi, On The Laplace Distribution of Firms
Growth Rates, L.E.M. Working Paper n. 2002-20, Scuola Superiore Sant'Anna
(2002).

G.Bottazzi, A Simple Micro-Model of Market Dynamics Part I:
The Large Market Deterministic Limit, L.E.M. Working Paper n.
2002-10, Scuola Superiore Sant'Anna (2002).

G.Bottazzi and G.Devetag, Coordination and Self-Organization
in Minority Games: Experimental Evidence, L.E.M. Working Paper n.
2002-09, Scuola Superiore Sant'Anna (2002).

G.Bottazzi, G.Dosi and G.Fagiolo, On the Ubiquitous Nature
of the Agglomeration Economies and their Diverse Determinants:
Some Notes, L.E.M. Working Paper n. 2001-10, Scuola Superiore Sant'Anna (2001).

G.Bottazzi, E.Cefis, G.Dosi, Corporate Growth and Industrial
Structure. Some Evidence from the Italian Manufacturing Industry,
L.E.M. Working Paper n. 2001-08, Scuola Superiore Sant'Anna (2001).

G.Bottazzi, G. Dosi and G. Rocchetti, Modes of Knowledge
Accumulation, Entry Regimes and Patterns of Industrial Evolution,
L.E.M. Working Paper n. 2001-06, Scuola Superiore Sant'Anna (2001)

G.Bottazzi, Firm Diversification and the Law of Proportionate Effect ,
L.E.M. Working Paper n. 2001-01, S.Anna School (2001).

G.Bottazzi, G.Dosi, M.Lippi, F.Pammolli, M.Riccaboni,
Processes of corporate growth in the evolution of an
innovation-driven industry: The case of pharmaceuticals, LEM
Workig Paper, Scuola Superiore Sant'Anna (2000)

G.Bottazzi, G, Devetag and G. Dosi, Adaptive Learning and
Emergent Coordination in Minority Games, L.E.M. Working Paper n.
1999-24, Scuola Superiore Sant'Anna (1999)


\section{Other Publications%
  \label{other-publications}%
}


\subsection{International journals and contributions to monographs%
  \label{international-journals-and-contributions-to-monographs}%
}

F. Cernuschi, N.Ludwig, P.Teruzzi and G.Bottazzi, A critical analysis
and possible modifications of two analytical model for defects
sizing using Video Pulse Thermography\textquotedbl{}, in D.Balageas, G.Busse and
G.M.Carlomagno Eds. \textquotedbl{}Quantitative Infrared Thermography\textquotedbl{}
Akademickie Centrum Graficzno-Marketingowe Lodart S.A. (Lodz
Polland 1999).

G. Bottazzi, G.Marchesini, G.P.Salam, M.Scorletti,, Small \$x\$ one
particle inclusive quantities in the CCFM approach, Journal of
High Energy Physics 12 (1998) 011

G. Bottazzi, G. Marchesini, G.P. Salam, M. Scorletti,, Structure
functions and angular ordering at small \$x\$,.  Nuclear Physics B
505 (1997) p. 366.

G. Bottazzi, G. Marchesini and M. Scorletti,, Final State Properties
in QCD Branching Processes at Small \$x\$, Nuclear Physics B
(Proc. Suppl.) 51C (1996) 250.


\subsection{Proceedings%
  \label{proceedings}%
}

G. Bottazzi, G. Marchesini and M. Scorletti,, Coherence and Final
States in DIS at Small \$x\$,.  in \textquotedbl{}Deep Inelastic Scattering and
Related Phenomena\textquotedbl{}, Proceedings of the DIS96 Conference, G.
D'Agostini and A. Nigro editors, World Scientific (Singapore
1997).

G. Bottazzi,, Angular ordering and structure function at small \$x\$,
Proceedings of the 5th Topical Seminar on 'The Irresistible Rise
of the Standard Model', San Miniato, Italy, 21-25 April 1997.


\subsection{Working Papers%
  \label{id1}%
}

G. Bottazzi, F.Cernuschi, N.Ludwig, P.Teruzzi, A critical analysis
and possible modifications of two analytical models for defect
sizing using Video Pulsed Thermography, QUIRT 98, Eurotherm
Seminar n.90, Polonia, 1998.

G. Bottazzi, Structure Function and final states with CCFM,
International Workshop on 'Low \$x\$ Physics at HERA', DESY
Zheuten (Berlin), Germany. June 1998.


\section{Software developed for research and teaching purposes%
  \label{software-developed-for-research-and-teaching-purposes}%
}

The following software is free and is distributed under the terms of
the GNU Public License at \url{http://cafim.sssup.it/~giulio/}

\begin{description}
\item[{gbutils}] 
A set of command line statistical utilities. Parametric and
non-parametric models together with various tests are
implemented. All the programs are written in C. Their output is
designed in a format suitable to be sent to popular plotting
programs. Since 2016 they are part of the scientific collection of
Debian/Ubuntu.

\item[{subbotools}] 
A set of programs written to deal with the Subbotin and antisymmetric
Subbotin distributions. Programs for maximum likelihood estimation of
parameters are included, together with generators of random
numbers. All the programs are written in C.

\item[{YAFiMM}] 
This package contains programs providing a simulation framework for
the analysis of a specific financial market model. For the
description of the model, see the Working Paper A Simple Micro-Model
of Market Dynamics Part I: The \textquotedbl{}Homogenous Agents\textquotedbl{} Deterministic
Limit in the WP page. All the programs are written in C++.

\item[{bose}] 
This package contains various programs related to the Bottazzi-Secchi
model of industry dynamics described in Bottazzi G. and Secchi A. On
the Laplace Distribution of Firms Growth Rates (see the WP
page). Some of the programs are written in Python, some in C.

\item[{NAIF}] 
NAIF (Numerical Analysis of Investment Functions) is a program to for
simulating the trading activities of an heterogeneous group og agents
in a simple two-good economy.

\item[{bdr}] 
This package contains the file needed to perform simulations of the
model of industrial dynamics described in G.Bottazzi, G.Dosi and
G.Rocchetti Modes of Knowledge Accumulation, Entry Regimes and
Patterns of Industrial Evolution Industrial and Corporate Change 10,3
(2001). All the programs are written in Python and require various
packages and modules that are freely and publicly available.

\end{description}

\emph{Autorizzo il trattamento dei miei dati personali ai sensi del Dlgs 196 del 30 giugno 2003}

\end{document}
